% status: FULL
\chapter{Numerics: Floating Point and Quantization Formats}\label{app:numerics}
\begin{ChapterIntent}
Ask which values a format represents, how arithmetic rounds, and which
equivalence the application requires. These questions precede bit counts.
Chapters~\ref{ch:weight-quantization}--\ref{ch:kv-quantization} provide the detailed low-precision formats and their
measured trade-offs; this chapter establishes the numerical vocabulary.
\end{ChapterIntent}

\section{Start with the gaps between numbers}
A stored floating-point value is not an arbitrary real number. At a
given magnitude, only a finite set of values is available. An arithmetic
result between two representable values must be rounded. A useful way
to understand precision is therefore to look at the gap between nearby
values, not only at the largest value a format can store.

Consider an illustrative binary format with three significant bits,
including the leading bit. Between one and two its values are
$1,\ 1.25,\ 1.5,\ 1.75$. At two, the exponent increases and the
spacing becomes one half. Adding a sufficiently small increment to
a large number can leave the stored value unchanged. No error message
is required; rounding is part of the arithmetic contract.

\EngineFigure{foundation-number-spacing}{An illustrative three-significant-bit
binary format. Equal relative precision produces larger absolute gaps at
larger magnitudes. The marked values are examples, not an FP8 encoding
table. Rounding maps an exact result to one representable point.}
{fig:foundation-number-spacing}

More exponent bits extend the available range. More significand bits
make values denser within a range. A format with fewer total bits can
trade one property for the other. This is why a datatype name must
identify a specific encoding, and why two low-precision formats with
the same bit count can behave differently.

Subnormal values, signed zero, infinity and NaN add boundary behaviour
outside the normal-number expression developed next. Some low-precision
encodings omit or reinterpret special values. A kernel must follow its
actual format, not a generic picture of floating point. The later
low-precision chapter reads concrete encodings and their scaling rules.

\section{Finite storage changes arithmetic}
A normal binary floating-point value has the schematic form
\[
 x=(-1)^s\,2^e(1+f),
\]
where $s$ is its sign, $e$ its unbiased exponent and $f$ the stored fraction
interpreted in $[0,1)$. This does not describe zeros, subnormals, infinities
or NaNs. Their encodings are format-specific.

Range and precision solve different problems. A larger range can avoid
overflow without resolving small changes near a large value. Flushing
subnormals matters near zero. Input, multiplication, accumulation and
output formats can differ within one operator; an FP16 label alone leaves
those questions unanswered.

\section{Separate storage, multiplication and accumulation}
A matrix product combines many products into each output element.
The weights can be stored in one format, multiplication can use another
representation and the accumulator can use a wider format. Saying
that a model is sixteen-bit does not fully describe that computation.

Imagine adding a thousand small positive products to a partial sum.
If the partial sum is repeatedly rounded in a coarse format, individual
increments may disappear. A wider accumulator can preserve more of
those increments, even when inputs remain low precision. The final
result may still be cast back to a smaller output type.

This does not make the operation exact. Input rounding occurred before
accumulation, and the accumulator has finite precision too. State what
each stage does. A comparison against a higher-precision reference
then has an interpretable meaning: it measures the combined effect
of the chosen numerical path.

Operation order also matters. Summing a row sequentially, in a tree,
or through several vector accumulators groups additions differently.
All can be mathematically valid and produce slightly different
floating-point results. Reproducibility requires a declared reduction
policy or a tolerance, not only a random seed.

\section{Stable softmax without overflow}
For scores $[1000,999]$, direct exponentiation can overflow. Subtract the
largest score:
\[
 \frac{[e^{1000},e^{999}]}{e^{1000}+e^{999}}
 =\frac{[1,e^{-1}]}{1+e^{-1}}
 \approx[0.731059,0.268941].
\]
Equality here is a real-arithmetic identity; the decimal values are rounded.
Adding the same finite constant to every score does not change exact
probabilities. A masked position contributes no mass. An entirely masked
row has no distribution under this normalization and needs an explicit
policy; blindly evaluating its maximum is not a policy.

A row summing to one can still be wrong. Compare its entries and downstream
weighted output, not merely normalization. Chapter~\ref{ch:flashattention}
uses this contract while changing the evaluation order.

\section{Specify the rounding rule}
Consider illustrative symmetric quantization with scale $s_q=0.5$:
\[
 q=\operatorname{clip}(\operatorname{round}(x/s_q),-7,7),
 \qquad \widehat{x}=s_q q.
\]
For $x=[-1.1,0.2,1.6]$, nearest rounding gives codes $[-2,0,3]$
and reconstruction $[-1,0,1.5]$. Absolute errors are $[0.1,0.2,0.1]$.
Without clipping, nearest rounding bounds error by $s_q/2$.
Clipping removes that bound: $x=10$ reconstructs as $3.5$, with error $6.5$.

At $x=1.25$ the quotient is $2.5$. Ties-to-even gives code 2;
ties-away-from-zero gives 3. The word round is therefore insufficient to
reproduce a conversion. The executable example uses ties-to-even and
tests this boundary. A production encoding must specify saturation,
nonfinite handling and scale computation too.

\section{Quantization is an encoding and a decoder}
Return to the scale-$0.5$ example. The code 3 reconstructs to 1.5
only because that scale is known. If another block uses scale 2,
the same code reconstructs to 6. A quantized weight therefore consists
of a code and the metadata needed to interpret it.

Run the committed arithmetic:
\begin{verbatim}
python3 code/foundations/lesson.py numerics
\end{verbatim}
The program prints codes $[-2,0,3]$, their reconstructions and the
block size, whose 18 bytes we account for in the next section.
Predict the output before running it. Then change
the illustrative input 1.6 to 10 in a scratch copy. The scale is
unchanged, so the code saturates at 7. This experiment separates
rounding error from clipping error with one visible number.

How should a scale be chosen? A simple symmetric choice covers the
largest absolute value in a group. An outlier then makes the step
coarser for every other element in that group. Smaller groups can
reduce this effect but require more scales. More elaborate methods
use calibration data or different error objectives; they are
introduced only after this basic encoding is clear.

Do not describe an integer code as a truncated float. Quantization
usually rescales and rounds the real value into a finite codebook.
Dequantization maps that code back to an approximation. Bit casting
would reinterpret an encoding without that numerical transformation.
The two operations are not interchangeable.

\section{Metadata is storage}
A hypothetical group of 32 four-bit codes and one 16-bit scale occupies
$32\cdot4+16=144$ bits: 18 bytes, or 4.5 bits per value.
This is an illustrative encoding, not every format named four-bit.
Zero points, secondary scales, padding and outlier storage change the count.
Smaller groups can improve local scaling while increasing metadata cost.

There is a computational bill as well: load and decode metadata, consume
packed codes, and accumulate products. Reduced payload does not guarantee
a faster unsupported kernel. Separate storage savings from measured runtime.

\section{Choose the equivalence}
Floating-point addition is not associative. In an illustrative decimal
system retaining three significant digits, adding $1000$, $1$, $-1000$
left-to-right loses the $1$; cancelling the large terms first retains it.
Binary formats exhibit the same phenomenon at their own spacing boundaries.

A numerical test can require
\[
 |\widehat y-y|\leq\epsilon_{\rm abs}+\epsilon_{\rm rel}|y|.
\]
Declare both tolerances; the absolute term matters near zero. If two logits
differ by a margin $\delta$ and each can move by at most $\epsilon$, their
ordering is guaranteed only when $\delta>2\epsilon$. Small tensor error,
the same greedy token and task quality are thus different tests.

A seed controls random draws, not rounding or reduction order. Quantization
also changes the model's numerical function. Report the intended contract
before deciding that a speedup preserves correctness; Chapter
\ref{ch:fast-wrong-answers} develops this hierarchy.

\section{Choose a test that answers your question}
Suppose reference logits are $[1.000,0.999]$. Moving the first down by
0.001 and the second up by 0.001 swaps the greedy winner. Each
coordinate changed only slightly, but generation can now follow a
different history. Once histories differ, later logits are evaluated
on different inputs, so they are no longer a direct kernel comparison.

A numerical test should therefore keep inputs fixed while comparing
operators. A generation test checks the accumulated effect on a
sequence. A task evaluation checks usefulness for a particular
workload. All three are valuable; none replaces the others.
The statement that a kernel passes a tolerance test does not establish
that a quantized model preserves task quality.

Relative error alone behaves badly near zero. If the true value is
zero, dividing by its magnitude is undefined. The combined absolute
and relative rule given earlier supplies a floor. Choose those
tolerances from the precision, operation and application, then record
them before judging the result. Increasing a tolerance only after
seeing a failure needs a reason, not an optimistic label.

\section{Use adversarially small examples}
Large random tests are useful, but a tiny structured example often
explains a failure faster. Use all-equal softmax scores to test
normalization, very large scores to test overflow handling, halfway
quantization values to test ties, out-of-range values to test clipping,
and cancellation between positive and negative terms to expose
reduction sensitivity.

Keep these examples when an optimized implementation arrives.
The optimized path might pack values, fuse operators or use a different
accumulation order. The examples still specify what it must preserve
and where approximation is allowed. An optimization is easier to
review when the expected numerical contract predates it.

At this point, fewer bits should no longer sound like a free speed
switch. They change storage, arithmetic support, conversion work and
error. The later quantization chapters will measure those trade-offs
on real implementations. Here we have built the vocabulary and the
small tests needed to read those measurements critically.
\section{Worked problems}
\subsection*{1. Rounding a tie}
\textbf{Problem.} With scale $0.5$, quantize 1.25 using ties-to-even.
\textbf{Solution.} The unrounded code is 2.5. The even code is 2,
which reconstructs to 1.0. Ties-away-from-zero would instead produce 1.5.
\subsection*{2. Preserve the winner}
\textbf{Problem.} The top-two logit margin is 0.03 and every logit changes
by at most 0.01. Can the winner swap?
\textbf{Solution.} No. The margin can shrink by at most 0.02 and remains
positive. At a margin of exactly 0.02 a tie could occur, so the strict
inequality matters.
\subsection*{3. Distinguish two errors}
\textbf{Problem.} Why does the half-step bound not cover the input 10
with scale 0.5 and codes limited to $[-7,7]$?
\textbf{Solution.} The nearest unbounded code would be 20, but clipping
forces code 7 and reconstruction 3.5. The 6.5 error comes mainly from
clipping, not from rounding to a neighbouring representable code.
\subsection*{4. Decode the scale bytes}
\textbf{Problem.} Interpret little-endian bytes \texttt{00 3E}
as FP16. Why would interpreting the same integer as BF16 be wrong?
\textbf{Solution.} Integer \texttt{0x3E00} has FP16 fields
$s=0,E=15,F=512$, so its value is $2^0(1+512/1024)=1.5$.
BF16 uses different field boundaries and bias. The byte count alone
does not identify the numerical format; BF16 1.5 is \texttt{0x3FC0}.
\subsection*{5. The stopped accumulator}
\textbf{Problem.} Why does adding one repeatedly in FP16 stop at 2,048,
and why does a final FP16 cast of the FP32 result still preserve 4,096?
\textbf{Solution.} At 2,048, spacing is two. The exact 2,049 is a tie,
resolved to the even retained significand at 2,048. Repeating presents
the same tie. FP32 keeps all 4,096 integer partial sums exact; 4,096
itself is a power of two representable in FP16. The lost information
depends on intermediate rounding, not only output type.
\subsection*{6. Bound a projection error}
\textbf{Problem.} For weights $[-1.1,0.2,1.6]$, reconstruction
$[-1,0,1.5]$ and input $[1,2,-1]$, calculate the output error
and its coordinatewise absolute bound.
\textbf{Solution.} The two dot products are $-2.3$ and $-2.5$.
Absolute error is 0.2. The triangle-inequality bound is
$0.1+0.2(2)+0.1=0.6$. Taking absolute values removes cancellation,
so the bound is valid but not attained here. This calculation
excludes arithmetic rounding after reconstruction.

